\documentclass[12pt]{article}
\usepackage{amsmath,amssymb,amsthm}
\usepackage[margin=1in]{geometry}

\newtheorem{theorem}{Theorem}
\newtheorem{lemma}[theorem]{Lemma}

\title{Problem 229: Four Representations Using Squares}
\author{Project Euler Solution}
\date{}

\begin{document}
\maketitle

\section{Problem Statement}
Count integers $n \leq N = 2 \times 10^9$ simultaneously representable as
\begin{align*}
n &= a^2 + b^2, \\
n &= a^2 + 7c^2, \\
n &= a^2 + 11d^2, \\
n &= a^2 + 13e^2,
\end{align*}
with non-negative integer variables (possibly different $a$ in each form).

\section{Mathematical Foundation}

\begin{theorem}[Sum of Two Squares]
A positive integer $n$ is a sum of two squares iff every prime $p \equiv 3 \pmod{4}$ appears to an even power in the factorization of $n$.
\end{theorem}
\begin{proof}
($\Rightarrow$) If $p \equiv 3 \pmod{4}$ and $p \mid a^2 + b^2$ with $p \nmid b$, then $(ab^{-1})^2 \equiv -1 \pmod{p}$, contradicting $\left(\frac{-1}{p}\right) = -1$. So $p \mid \gcd(a,b)$ and $p^2 \mid n$; by induction, $p$ appears to even power.

($\Leftarrow$) Primes $p \equiv 1 \pmod{4}$ split in $\mathbb{Z}[i]$ as $p = \pi\bar{\pi}$ with $N(\pi) = p$, so $p$ is a sum of two squares. We have $2 = 1^2 + 1^2$. By the Brahmagupta--Fibonacci identity, the product of sums of two squares is a sum of two squares. Primes $p \equiv 3 \pmod{4}$ to even power give $p^{2k} = (p^k)^2 + 0^2$.
\end{proof}

\begin{theorem}[Sieve Correctness]
For each $k \in \{1, 7, 11, 13\}$, the set of representable $n \leq N$ is exactly $\{a^2 + kb^2 : a, b \geq 0,\, a^2 + kb^2 \leq N\}$.
\end{theorem}
\begin{proof}
By definition, exhaustive enumeration of all valid $(a,b)$ pairs.
\end{proof}

\begin{lemma}[Bit Array Intersection]
The answer equals $\mathrm{popcount}(B_1 \wedge B_7 \wedge B_{11} \wedge B_{13})$, where $B_k[n] = 1$ iff $n$ is representable as $a^2 + kb^2$.
\end{lemma}
\begin{proof}
Bitwise AND computes set intersection; popcount gives cardinality.
\end{proof}

\begin{lemma}[Pair Count]
The number of pairs $(a,b)$ with $a^2 + kb^2 \leq N$ is $\Theta(N/\sqrt{k})$.
\end{lemma}
\begin{proof}
The count approximates $\int_0^{\sqrt{N/k}} \sqrt{N - kt^2}\,dt = \frac{\pi N}{4\sqrt{k}}$.
\end{proof}

\section{Algorithm}
\begin{enumerate}
\item Divide $[0, N]$ into blocks of size $B = 10^7$.
\item For each block, generate all representable numbers in that range for all four forms using bit arrays.
\item AND the four bit arrays; accumulate the popcount.
\end{enumerate}

\section{Complexity Analysis}
\begin{itemize}
\item \textbf{Time:} $O(N)$ total across all forms.
\item \textbf{Space:} $O(B/8)$ per bit array, $\sim 5$\,MB total with $B = 10^7$.
\end{itemize}

\section{Answer}

\[
\boxed{11325263}
\]

\end{document}
